\section{降低未冻结站点之间的宽度}\label{subsec:core-width}

集中操作之后，有些站点在本轮余下的过程中保持冻结，其余站点构成\emph{核心}。我们先降低核心站点间的宽度，再恢复有界度数，并把所有站点交给下一轮。每个核心位置实例承载一个张量。每个冻结位置实例持有一个张量列表，并记录每个张量冻结时的计费量。所有外腿都接在冻结张量上。冻结站点间的宿主邻接关系不承载任何连线，可以删除。

设 \(X\) 为所有块中核心站点与核心站点间连线的宽度之和，每条线只计一次。设 \(Z\) 为冻结站点的总计费量。核心与冻结站点之间的连线称为\emph{系线}，其总宽度至多为 \(Z\)。装包时，将每个核心位置实例的负载取为其元数，将每个冻结位置实例的负载取为其计费量之和乘以属于 \((0,1]\) 的固定二进因子。用 \(B\) 表示单个负载的上界。所有负载之和满足 \(\mathcal M\le 2X+2Z\)。存在一个共同的整数 \(u\le B\)，使所有内部宽度都是 \(u\) 的倍数。

宿主图记录所有块中可能出现的接触。用 \(\Delta\ge3\) 表示其核心度数的上界，用 \(\nabla\ge\Delta\) 表示整个宿主图的最大度数的上界。核心站点在不同块中可能与不同的冻结站点接触，因此即使核心的度数为三，仍可能有许多潜在系线。下面的宽度估计与潜在系线的条数无关。

\begin{lemma}[核心宽度的缩减]\label{lem:core-width} 固定输入核心度数的一个上界，并设定理~\ref{thm:sort} 的假设成立。存在一个固定的分组及收缩步骤序列，称为一\emph{拍}。给定余量 \(\sigma>0\)，每个求和项都满足 \[
 X_{\rm out}\le \tfrac12 X+C(Z+\sigma),\qquad
 \Delta_{\rm out}\le3,\qquad Z_{\rm out}=Z.
\] 这一恒等式分别保留每条外腿。最大负载、线束单位的倒数和 \(\log(2\nabla)\) 的增长倍数均有固定上界。这些因子及 \(C\) 均与输入的全图最大度数无关。

在宿主图有 \(H\) 个位置、块数 \(N_b\) 为二的幂的一步中，选取二的幂 \(r\mid N_b\)。分组后剩下 \(N_b/r\) 个块，至多使用 \(crH\) 个位置。只须在每步取 \(r\ge c_*\)、\(c_*HN_bB/r\le\sigma\)，并使初始单位能被各细分因子的乘积整除。位置常数及余量常数 \(c,c_*\) 可以依赖于相应步骤的度数上界。\end{lemma}

一拍由若干节省连线的子轮和一次重置组成。在每个子轮中，允许可能的核心接触增多，以换取核心宽度的缩减。重置时，将余下的连线经过一个度数为三的核心布设。我们将选择子轮数，使各子轮的节省合起来足以抵付这次布线的开销。

\subsection{一个接线恒等式} 将六方按图~\ref{fig:two-layouts} 排成三角阵列。若两方同在一行或一列，则称这一对为\emph{允许对}。所需的相等连线 \(AB\) 和 \(CD\) 的两端既不同行也不同列。我们将其乘积表示为若干只在允许对上布线的网络之和，随后按行或按列收缩整个划分。这样，同一个求和项有两种布局可选。

\begin{figure}[htb]\centering \begin{tikzpicture}[x=1.06cm,y=1.32cm,font=\small,
 party/.style={circle,inner sep=1pt,minimum size=4mm},
 part/.style={rounded corners=2.8pt,line width=.65pt},
 demand/.style={draw=ink,densely dashed,thick}]
\foreach \panel/\shift/\heading in {0/0/需求连线,1/4.2/按行划分,2/8.4/按列划分}{
 \begin{scope}[xshift=\shift cm]
 \node[font=\sffamily\small] at (1.14,1.96) {\heading};
 \ifnum\panel=1
  \draw[part,draw=sea,fill=sea!8] (-.33,1.15) rectangle (2.61,1.65);
  \draw[part,draw=sea,fill=sea!8] (-.33,.45) rectangle (1.47,.95);
  \draw[part,draw=sea,fill=sea!8] (-.33,-.25) rectangle (.33,.25);
 \fi
 \ifnum\panel=2
  \draw[part,draw=rust,fill=rust!8] (-.33,-.25) rectangle (.33,1.65);
  \draw[part,draw=rust,fill=rust!8] (.81,.45) rectangle (1.47,1.65);
  \draw[part,draw=rust,fill=rust!8] (1.95,1.15) rectangle (2.61,1.65);
 \fi
 \ifnum\panel=0
  \draw[demand,shorten <=2.6mm,shorten >=2.6mm] (0,1.4)--(1.14,.7);
  \draw[demand,shorten <=2.6mm,shorten >=2.6mm] (2.28,1.4)
    .. controls (2.64,-.14) and (1.26,-.23) .. (0,0);
 \fi
 \foreach \x/\y/\name in {0/1.4/C,1.14/1.4/M,2.28/1.4/A,0/.7/N',1.14/.7/D,0/0/B}{
  \node[party] at (\x,\y) {\(\name\)};
 }
 \end{scope}}
\end{tikzpicture} \caption{例如，\(C\) 和 \(M\) 之间的一条相等连线在按行收缩时消失，在按列收缩时保留。每条允许连线在其中一种布局下都会消失。对于一个求和项中该恒等式的所有副本，统一采用两种布局中代价较低的一种，所保留的宽度至多为替代连线总宽度的一半。}\label{fig:two-layouts} \end{figure}

一条 \(s\) 比特连线是两个 \(s\) 比特指标上的相等张量。对于 \(k\) 上一个尺度为一的普通价格，令 \(m\) 表示允许对上最大的价格距离。在允许图的一棵生成树上，令每个顶点处的关联比特相等，就能提供一个六方共享的比特。因此 \(1\le5m\)。两个所需距离都至多为 \(2m\)。若 \(d(C,D)=2m\)，则路径 \(C,M,D\) 和 \(C,N',D\) 使 \(M\) 和 \(N'\) 在一个 CAT(0) 扩张中成为同一个中点。若还满足 \(d(A,B)=2m\)，那么 \(C\) 与 \(M=N'\) 都是 \(A,B\) 的中点。于是 \(C=M\)，这与 \(d(C,M)=m>0\) 矛盾。由尺度为一的价格的紧致性，存在统一的 \(\lambda_*<2\)，满足 \[
 d(A,B)+d(C,D)\le2\lambda_*m .
\]

取 \(\lambda_*<\lambda_0<2\)。我们断言，对固定的 \(\eta'>0\) 和任意大的 \(b\)，两条宽度为 \(b\) 的需求连线可以写成至多 \(2^{\eta'b}\) 个网络的普通和，每个网络的允许连线总宽度至多为 \(2\lambda_0b\)。事实上，选充分大的 \(h_0\) 使 \(\lfloor2\lambda_0h_0\rfloor>2\lambda_*h_0\)。作为供给，在每个允许对上取一条宽度为 \(\lfloor2\lambda_0h_0\rfloor\) 的连线，再对所有允许对取直和。在每个谱点，这一供给的价格都超过两条宽度为 \(h_0\) 的需求连线的价格。由严格谱比较，对 \(b=h_0r\)，宽度为 \(b\) 的需求是供给的 \(r\) 次幂的有界多个副本之直和的一个限制。展开直和并施以局部限制映射，即得所称网络。记允许对的集合为 \(E_{\rm perm}\)，一个求和项的指标需要 \(r\log_2|E_{\rm perm}|+O_{h_0}(1)\) 比特。先取充分大的 \(h_0\)，再取充分大的 \(r\)，使上述比特数不超过 \(\eta'b\)。\(k\) 上的该恒等式可扩张至任意扩域 \(K/k\)。为构造子轮，我们需要即使每条线承载同一比特的许多副本，选择的数目仍然固定。

\begin{lemma}[宽度的固定选择]\label{lem:powered-gadget} 存在常数 $b\in\mathbb N$ 和 $\lambda<2$，使得对每个 $z\ge1$，两条宽度为 $bz$ 的需求连线等于同一组六方上有界多个网络的和。每个网络在允许边上至多使用 $2\lambda b$ 个线束，每束含 $z$ 比特。局部张量由 $z$、所在的一方及求和项指标指定。指标集及各宽度的选取（以 $z$ 为单位）可以独立于 $z$ 固定。\end{lemma}

\begin{proof} 将上述有限恒等式取 $z$ 次张量幂。起初，每一项由含 $z$ 个选择符的序列标记，每个选择符至多占 $\lceil\eta'b\rceil$ 比特。在每条允许边上，把比特总数向上取整至 $z$ 的倍数。按取整后以 $z$ 为单位的各边宽度向量对序列分组，允许空组。该向量取值于固定有限集，每条允许边为取整多付至多一个单位。每一组中，取连通的允许图的一棵生成树，沿每条树边传送完整序列，每边额外使用至多 $\lceil\eta'b\rceil$ 个单位。局部表检验所有关联树边上的序列是否相同且有效、是否属于该组，并检验填充位置是否为零。它使用该序列将每条原边的位置划为规定的片段，再计算 $z$ 张局部表的乘积。由连通性，一旦指定序列，它在各树边上的编码就是唯一的。因此，对树变量求和恰好给出该组的和。以 $z$ 为单位的总数至多为 $2\lambda_0 b+|E_{\rm perm}|+5\lceil\eta'b\rceil$。先取充分小的 $\eta'$，再取充分大的 $b$，就能使该量至多为 $2\lambda b$，其中 $\lambda<2$。\end{proof}

\subsection{子轮} 在一个子轮中，各个位置的位置实例被分别装包。来自不同包的两个出向单位可以分别充当 \(AB\) 和 \(CD\)。先安排足够多这样的配对，并保证单次收缩所合并的包数有界。

\paragraph{平衡划分。} 对每一组 \(r\) 个块指标，按共同顺序取二进区间。给每条宿主边 \(e\) 指定传播距离 \[
 s_e=64\lceil\log_2(2x_e)\rceil,\qquad
 x_e=\begin{cases}\Delta&\text{在核心边上},\\
                     \nabla&\text{在系线上}.\end{cases}
\] 在每个位置标记负载大于 \(AB\) 的所有区间，其中 \(2\le A\le r\) 为二的幂。标记意味着必须分裂。将每个标记传播给同一位置的父区间，并且沿每条关联边 \(e\)，传播给另一端高 \(s_e\) 层的祖先区间，只要该祖先存在。叶区间称为聚落，其负载至多为 \(AB\)，长度至少为 \(A\)。由于较小聚落的父区间所发出的传播，边 \(e\) 两端相交的聚落长度之比至多为 \(2^{s_e}\)。

在每组内部对指标采用一个公共置换，可使聚落总数满足 \begin{equation}
 Q\le C_1\left( HN_b/r+\mathcal M/(AB)\right)
 \label{eq:core-colonies}
\end{equation} 其中 \(C_1\) 为绝对常数。为看出这一点，在一个二进制基下把指标等同于 \(\mathbb F_r\)，再乘以均匀随机选取的非零元。两个相异指标落入同一深度 \(j\) 区间的概率至多为 \(2^{-j}\)。对一组中的一个源位置，记 \(x_i\le1/A\) 为各负载除以 \(AB\) 之后的值，令 \(x=\sum_i x_i\)。初始即带标记的区间内，不同归一化负载的有序交叉乘积之和至少为 \(1/2\)。因此，深度 \(j\) 处的初始标记数的期望至多为 \(2x^2 2^{-j}\)。在每个顶点，对关联边求 \(2^{-s_e/16}\) 之和至多为 \(1/2\)。对所得路径权重求和可知，总传播距离至多为 \(s\) 时，可达位置至多有 \(2\cdot2^{s/16}\) 个。固定总向上位移 \(s\)，允许在同一位置向父区间迈步。那么在每个可达位置，深度 \(j\) 处的源标记到达的不同祖先至多有 \(2^{j-s}\) 个。在该位置对 \(j\ge s\) 求和，期望数至多为 \[
 \sum_{j\ge s}\min\{2^{j-s},2x^2 2^{-j}\}=O(x2^{-s/2}).
\] 再对距离、源和组求和，标记总数的期望至多为 \(O(\mathcal M/(AB))\)。每个二叉划分的叶数比带标记的内部节点数多一。选择使实际数目最小的乘子即可。

\paragraph{节省连线的一次子轮。} 为所有核心边定向。用 \(\mathbb F_2\) 上的一个仿射泛函作用于块在组内的指标，将每个块着为粗或细两色之一。同一块的所有位置颜色相同。在同一个位置装包时，取各张量的张量积，同时保留连线。随后合并产生的自环收缩到局部表内。在每个核心聚落，把所有粗张量装在一个包内。其余依次装成细包，负载达到 \(B\) 时即封包。细包共 \(O(A)\) 个，每包负载至多为 \(2B\)。在冻结聚落，把各列表拼接起来。记 \(P=2^{64\lceil\log_2(2\Delta)\rceil}\) 及 \(K_1=\Delta P\)。由颜色的两两独立性，每个聚落中粗、细两类出向宽度之差的绝对值的期望至多为 \((B\sum_i a_i)^{1/2}\)，其中 \(a_i\) 是其第 \(i\) 个位置实例的出向核心宽度。求和可知，有一种着色使总不平衡量至多为 \(\sqrt{BQX}\)。

对每个源聚落，按目的包把粗出向单位分箱。由核心边两端的长度比可知，箱数至多为 \(K_1\)。把每个细包的所有出向单位视作一件物品，将整件物品匹配至任何一个剩余容量足够的箱中。首次匹配失败时，每个箱的剩余容量都小于 \(2B\)，故两侧总量中较小者的未匹配部分至多为 \(2BK_1\)。因此，按单侧计的匹配宽度满足 \[
 I\ge \tfrac12X-\tfrac12\sqrt{BQX}-2BK_1Q.
\] 聚落数上界~\eqref{eq:core-colonies} 与取 \(c_*\ge A\) 时的余量条件给出 \(BQ=O((X+Z+\sigma)/A)\)。取 \(A\) 为 \(1+K_1\) 的充分大的二的幂倍数，必要时取整。于是 \(I\ge X/4-O(Z+\sigma)\)，隐含常数为绝对常数。

在每个箱内，将已匹配的细单位与互异的粗单位配对。应用引理~\ref{lem:powered-gadget}，取 \(b_*=b\) 和 \(z=u/b_*\)；由细分假设后者为整数。在每一对中，粗线充当 \(AB\)，其源为 \(A\)，细线充当 \(CD\)，其源为 \(C\)。两个源属于同一聚落。由于一条线始终留在一个块内，且整个块颜色相同，两个目的端 \(B,D\) 分别属于粗包和细包。为每一对引入新的辅助方 \(M,N'\)。遍及所有配对，两条需求连线的宽度总和为 \(2I\)，每个乘积型求和项中替代线的总宽度至多为 \(2\lambda I\)。

在其中一个求和项内，对所有配对按行收缩后保留的宽度求和，再对按列收缩的情形作同样计算。每条替代线在两种布局之一中成为内部线，因此这两个总量之和至多为 \(2\lambda I\)。在所有配对中统一执行代价较小的那种收缩。进一步认同顶点只会使保留的宽度减小。未匹配连线的贡献至多为 \(X-2I\)，故 \[
 X'\le X-(2-\lambda)I\le\rho X+C(Z+\sigma),\qquad
 \rho=1-(2-\lambda)/4<1 .
\]

以整件为单位匹配每个细包，可以同时控制负载和度数的增长。被合并的旧包构成星形图。细包作为源时，在行布局中与自身所在聚落的粗包合并，在列布局中与其唯一选定的粗目的包合并。作为目的时，它仅与一个辅助方合并。因此每个粗中心至多吸收 \(O((1+K_1)A)\) 个细包。每个辅助方都并入一个需求端点，可以将其新增关联宽度按匹配宽度的常数倍计入该端点。可用一个二的幂次负载增长因子 \(g\)，同时控制所得核心负载和冻结列表的缩放负载，且该因子在 \(\Delta\) 的多项式以内。新单位为 \(u/b_*\)。

在每个可能的二进区间中预留包槽位。可能的核心接触只发生在同一聚落的包之间，或相交且长度比至多为 \(P\) 的相邻区间的包之间。对一个给定区间，在每个邻点有 \(O(P+\log(2P))\) 个可能区间。上述星形图及四个匹配端点在这张图中的直径都有界。取图的一个固定次幂，就得到度数至多为 \(C_2\Delta^{C_2}\) 的公共核心支撑图。系线只需经过一处具有自身长度比界的旧系线邻接，再经过有界多个核心邻接。因此全图最大度数至多为 \(C_2\nabla^{C_2}\)，其中绝对常数 \(C_2\) 必要时可增大。每个位置的候选区间少于 \(2r\) 个。这样便得到了所称的位置数界，并使各系线仍接至各自的冻结张量。

\subsection{重置核心度数}\label{subsec:reset} 以 \(H,N_b,\Delta,\nabla\) 表示重置前一刻的参数，并令 \(P=2^{64\lceil\log_2(2\Delta)\rceil}\)、\(K_1=\Delta P\)。重置用一个度数为三的宿主图替换核心。其主要开销是每个核心基区中的一个置换网络。固定一个规模为 \(c_0Yw_0\) 的三价网络，使得通过固定其标签，能够实现 \(Y\) 行（每行宽 \(w_0\)）的所有置换，其中 \(4K_1\le Y\) 且 \(c_0\ge1\)。称这一网络为\emph{转子}。

图~\ref{fig:reset} 给出两种路径。先根据平均负载，把每个位置的块划为称作\emph{基区}的区间，并在每对相邻核心基区之间预留容量。混排各块使大多数基区能符合这些预留容量。两个这样的基区之间的线，直接经过两个专用转子输出端之间的连接。与例外基区关联的每条核心线则沿二叉树走较长的路径。要将这些额外跳数计入预算，必须先控制例外基区的总负载。

\begin{figure}[htb]\centering \begin{tikzpicture}[x=1cm,y=1cm,
                        font=\footnotesize,term/.style={circle,draw,fill=white,inner sep=2pt}]
% pair of good bases
\foreach \d/\off in {1/-5,-1/5}{
  \begin{scope}[xshift=\off cm,xscale=\d]
    \draw[draw=black!60,rounded corners=6pt](-1.45,-1.5) rectangle (3.36,1.5);
    \node[font=\sffamily\small] at (1,1.8){良好基区};
    \node at (-.05,1.18){输入};
    \node at (2.6,1.18){输出};
    \foreach \y in {.75,0,-.75}{
      \draw (0,\y) -- (2.6,\y);
      \draw (0,\y) -- (-.65,\y) --(-.87,\y+.21);
      \draw (-.65,\y)--(-.87,\y-.21);
    }
    \draw[sea, very thick] (0,0)--(.8,0) (1.8,.75)--(2.6,.75);
    \draw[draw=ink,rounded corners,fill=ink!8] (.8,-1.26) rectangle (1.8,1.12);
    \draw[sea,very thick] (.8,0) -- (1.8,.75);
    \node[font=\scriptsize] at (1.3,-.59){转子};
    \foreach \y in {.75,0,-.75}{
       \node[term] at (0,\y){};
       \node[term] at (2.6,\y){};
    }
  \end{scope}}
\draw[sea,very thick] (-2.29,.75) -- node[below,text=black]{预留连线} (2.29,.75);

% link from bottom input on left to a terminal in a bad base
\draw[draw=black!60,rounded corners=5pt] (-2.15,-3.94) rectangle (0.2,-2.23);
\node[anchor=west,align=center] at (.42,-3.1) {不良基区的叶};
\node at (-.61,-2.7) {输入};
\draw (-1.41,-3.26) --(-1.68,-2.9);
\draw[rust,thick,densely dashed] (-5.11,-.75)--(-5.65,-.75)--(-5.87,-.96)
  -- (-5.87,-3.6)-- (-1.68,-3.6) -- (-1.41,-3.26)--(-.6,-3.26);
\node[term] at (-.6,-3.26){};
\node[align=center] at (-4,-2.91) {经接触树};
\end{tikzpicture} \caption{重置时的布线。在良好基区之间，连线的每个端点都通过一个转子抵达专用输出端（实线）。与不良基区的接触通过二叉接触树连接输入端（虚线）。旧张量及其系线留在输入端。图中只画出部分端口。}\label{fig:reset} \end{figure}

\paragraph{基区与容量。} 取二的幂 $F\ge Y^2w_0$，本步骤要求 $c_*\ge F$。对各位置在 \(N_b\) 个块上的平均负载 \(\mu_v\)，令 \[
 \nu_v=\max\left\{ FB/r,\ 
      \max_{\pi:y\to v} \mu_y\,2^{-\sum_{e\in\pi}s_e/2}\right\},
\] 允许长度为零的路径。与计数聚落时一样，从每个位置出发的各路径的衰减权重之和受一个绝对常数控制。因此 \(\sum_v\nu_v\le HFB/r+C\sum_v\mu_v\)。在位置 \(v\) 处分成等长基区，长度取满足 \(s_v\le FB/\nu_v\) 的最大二的幂，介于 \(F\) 与 \(r\) 之间。各次块指标的划分均与每 $r$ 个连续指标一组的分组相容。相邻基区的深度之差至多为 \(s_e/2\)。基区总数为 \(O((\mathcal M+\sigma)/(FB))\)。对相邻核心位置 \(v,w\) 上每对相交的基区，预留容量 \[
 \overline y_{vw}\min(s_v,s_w)+FB/K_1 ,
\] 其中 \(\overline y_{vw}\) 是按块平均的该接触的核心宽度。在任一基区，第一项之和至多为 \(s_v\mu_v\le FB\)，而预留项至多有 \(K_1\) 个，故该基区的总预留容量至多为 \(2FB\)。在置换各块之前，基区长度及容量已由均值固定。因而对每一组，给定一对基区位置的输出端将占用相同的宿主槽位。

\paragraph{混排。} 用同一个可逆二元矩阵置换所有块指标。在每个基区内，标记负载超过 \(2FB\) 的二进子区间（可能包括基区本身），按相同的边距离传播。叶区间负载至多为 \(2FB\)，长度至少为 \(F\)，并满足先前的长度比界（比较基区深度可知，标记区域上方的基区分裂已满足传播要求，再按聚落装包时的方法传播）。若基区被分裂，或与之关联的核心预留容量溢出，则称其为不良基区。对任意 \(\delta_*>0\)，按 \(\Delta,Y,\delta_*\) 选足够大的 \(F\)，即存在一个置换使 \begin{equation}
 \sum_{\text{不良基区}}\mathrm{load}\le\delta_*(\mathcal M+\sigma).
 \label{eq:core-badbases}
\end{equation} 为了使证明对全图最大度数一致，计算一个过载区间的传播范围。对有 \(n_v=N_b/s_v\) 个基区的某个给定位置，将原负载除以 \(FB\) 归一化，得到 \(x_i\le1/F\)。在低于基区 \(j\) 层的深度，设 \(a_J\) 为区间的归一化负载，\(\overline a\) 为其均值，它至多为 \(2^{-j}\)。量 \(V_j=\sum_J(a_J-\overline a)^2\) 和 \(U_j=\sum_J a_J^2-\sum_i x_i^2\) 各自均为过载区间数及其归一化负载总和的一半的上界。在随机可逆二元矩阵作用下，不同指标发生前缀碰撞的概率至多为区间数的倒数。因此 \[
 \mathbb E V_j\le\sum_i x_i^2\le n_v/F,\qquad
 \mathbb E U_j\le n_v2^{-j}.
\] 由于相邻基区的深度差至多为边距离的一半，深度 \(j\) 的标记只有沿满足 \(\sum s_e\le 2j\) 的路径才能迫使某个基区分裂。对该源标记，由二进区间的嵌套性，每个可达位置至多有一个基区受影响。因此由~\eqref{eq:core-colonies} 所用的路径计数，受影响数目至多为 \(2\cdot 2^{j/8}\)。在深度 \(j\) 处，若 \(1/F\le2^{-j}\) 则用 \(V_j\)，否则用 \(U_j\)。给这一估计量赋权 \(4\cdot 2^{\lceil j/8\rceil}\)，在 \(j=\log_2 F\) 处分开求和。期望计费量至多为 \(n_v\sum_{j\ge0}2^{j/8}\min\{F^{-1},2^{-j}\}=O(n_vF^{-7/8})\) 的一个常数倍。负载超过 \(2FB\) 的基区直接由深度零的估计支付。对核心容量溢出，在基区相交的深度，对接触宽度除以 \(FB\) 的结果应用方差估计量。预留容量比均值高出 \(FB/K_1\)，保证每次溢出至少贡献 \(1/K_1^2\)。将无权方差的期望对核心边求和，成本至多为 \(X/(F^2B)\)。每次溢出至多影响两个基区。其中负载超过 \(2FB\) 的已在深度零处计费，其余每个受影响基区的成本至多为 \(2FB\)。令 \(\Phi\) 为加权区间估计量之和再加上各接触方差之和的 \(K_1^2\) 倍。以上给出了不良基区总负载的确定性上界 $O(FB\Phi)$，并证明 \[
 FB\,\mathbb E\Phi = O(F^{-7/8}(\mathcal M+\sigma)+ K_1^2X/F).
\] 取估计量不超过其均值的一个矩阵，即得 \eqref{eq:core-badbases}。第~\ref{sec:assembly} 节通过计算碰撞得分找出这样的置换。

\paragraph{布线路径。} 混排后，将连续的 \(r\) 个块分成一组。给每个核心叶区间配置 \(Y\) 个输入端，把它的各旧张量逐个完整分配给当前负载最小的端口，使每个端口接收的负载至多为 \(2FB/Y+B\le 3FB/Y\)。所有系线都留在分配到的端口上，其另一端仍在拼接后的冻结列表中。在每个良好核心基区，转子将输入端连接至 \(Y\) 个输出端。对一项容量为 \(c\) 的预留，分配 \(\lceil cY/(8FB)\rceil\) 对专用输出端，每对均以直接连线相接。这样每个基区至多需要 \(Y/4+K_1\le Y\) 个端口。将两个良好基区之间一次接触的所有单位尽可能均匀地分到相应各对端口，输出负载至多为 \(8FB/Y+u\le9FB/Y\)。

在一个良好基区，记 \(m_{ij}\) 为从输入 \(i\) 路由到输出 \(j\) 的单位数。矩阵 \((\lceil m_{ij}/w_0\rceil)\) 的各行和及各列和均至多为 \[
 9FB/(Yuw_0)+Y\le10FB/(Yuw_0).
\] 通过添边把所有行和、列和都补到其中的最大值。由 Hall 定理，所得正则二部多重图可分解成完美匹配。每个匹配配置转子的一层，其中比特边取张量幂，使宽度达到 \(u\)。把未用的输入及其对应输出固定为零。在转子的每个位置，取各层表的张量积。这只用到通过固定转子标签获得的规定相等张量。在转子各行，平衡二叉树汇集其 \(w_0\) 个端口，每个单位经过 \(O(\log(2w_0))\) 跳。计入直接连线后，每个良好基区的开销为 \(O((c_0+\log(2w_0))FB)\)。

对一端不良的核心接触，使用接在相应输入端上的平衡二叉接触树。可能相邻的叶区间长度比至多为 \(P\)，每叶有 \(Y\) 个端口，故深度为 \(O(\log(2\Delta Y))\)。利用相等张量沿树传送有序列表。两端已分配的列表唯一确定所传比特的内部赋值。在 \eqref{eq:core-badbases} 中取充分小的 \(\delta_*\)，即足以支付该布线深度。因此，对某个绝对常数 \(A_0\)，有 \[
 X'\le A_0(1+c_0+\log(2w_0))(X+Z+\sigma).
\] 这里计入了转子、汇集树、直接连接以及接触树中的每一条核心连线。单位仍为 \(u\)。由端口与层的界，每个输入端、输出端、树顶点及转子顶点的负载都是 \(O(FB)\)。每个冻结叶区间中，缩放后的列表负载至多为 \(2FB\)。核心度数至多为三。转子顶点有此度数。在核心内，一个输入端至多与一个转子行（通过一条接往行树的边）及一棵接触树相连，一个输出端至多与一个转子行及一条直接连线相接。每棵二叉树度数至多为三。对系线，公共宿主图允许连接到长度比符合规定且可能相交的所有核心叶区间上的每个输入端。所得全图最大度数以 \(2\nabla\) 的某个固定次幂为界，其中的常数依赖于此次重置。在所有候选区间槽位上预留这些设施至多需要 \(c r H\) 个位置。

\subsection{支付重置的开销} 用 \[
 \kappa=A_0(1+c_0+\log(2w_0))
\] 表示重置的宽度乘数。取 \(t\) 为满足 \(\kappa\rho^t\le1/2\) 的最小正整数。反复使用子轮估计可知，在重置之前，核心宽度至多为 \(\rho^tX+C(Z+\sigma)/(1-\rho)\)。重置后，输入宽度 \(X\) 的系数至多为二分之一。其余开销为 \(Z+\sigma\) 的固定倍数，条件是对经过 \(t\) 个子轮后达到的度数，能选到满足 \(Y\ge4K_1\) 的转子。

起始核心度数有界时，递推 \(\Delta'\le C_2\Delta^{C_2}\) 给出 \[
 \log_2(4K_1)\le C^{t+1}
      \le C''(2+c_0+\log(2w_0))^{C'} .
\] 这里 \(C,C',C''\) 仅依赖于输入核心度数和各子轮常数。对 \(Y=2^y\)，由定理~\ref{thm:sort}，对每个 \(\zeta>0\) 可得 \(w_0=\lceil y^{a_\zeta}\rceil\) 及 \(c_0=O_\zeta(y^\zeta)\)。固定 \(\zeta<1/(2C')\)，再取充分大的 \(y\)。上面给出的多项式界至多为 \(y\)，故该转子可处理所需的核心度数。

子轮数 \(t\)、转子、\(F\) 及 \(\delta_*\) 仅取决于输入核心度数的上界和条件性的排序构造。这些参数控制最大负载的增长。全图最大度数决定系线传播的范围和要预留的宿主位置数。经过全部子轮，\(u\) 除以 \(b_*^t\)；各子轮中的多项式度数界再与重置合并，使 \(\log(2\nabla)\) 至多乘以一个固定因子。

引理~\ref{lem:core-width} 得证。
